% @Author: Jacem Chaieb
% @Date:   2015-07-26 13:42:06
% @Last Modified by:   Jacem Chaieb
% @Last Modified time: 2016-04-12 15:49:40

%%%%%%%%%%%%%%%%%%%%%%%%%%%%
% CHAPTER                  %
%%%%%%%%%%%%%%%%%%%%%%%%%%%%
\chapter*{Introduction}
\label{chap:general_intorduction}
\markboth{\MakeUppercase{Introduction}}{}%
\addcontentsline{toc}{chapter}{Introduction}%
The internet since its inception has been a path humans took to share data, share resources, and overall share knowledge, with the recent events that hit the world it has been proven that the internet 
can be used to share more than object raw data, but we can effectively share our lively-hood, with real time communication we can be with other people the other end of the pacific in milliseconds, but 
it the human aspect was always lacking, so the question that asks itself now, is how can we make the experience as seamless as real life ? an ambitious take on web communication, so my host company Nethermind, 
decided to take on this problem by providing what is known as a Metaverse, this project has as a goal providing means by which users can communicate in real time while being able to interact with each other 
thought avatar of their choice, a gamified replica of the real world in a virtual space.\\
in this context our task is to develop an instance of a general purpose metaverse, and an art specific art-verse, and to realise our project we will follow the following approach : \\
\begin{itemize}
    \item Chapter 1 :  in this chapter we will get to know the host company, the project scope and goals, and we will know more about the team's methodology and way of tackling various aspects of the project;
    \item Chapter 2 :  in the second chapter we will identify the requirements and needs that the system needs to fulfill, as well as the actors and target audience of said project;
    \item Chapter 3 :  the third chapter will be dedicated to the abstract conceptualization of the project, it is a theoretical and technical study of the project do-ability;
    \item Chapter 4 : the fourth and final chapter will be allocated to the realisation of the project discussing tools, methods and algorithms utilised to implement various aspects of the project;
\end{itemize}